\chapter{Méthodologie et organisation du projet}
\section{Démarche}
Ce chapitre rappelle la démarche méthodologique adoptée pour la conduite du projet, depuis le cadrage initial jusqu'à la mise en production progressive de la solution.

Une approche itérative et incrémentale, inspirée des méthodes agiles (Scrum), a été retenue afin de livrer rapidement une première version fonctionnelle du système, puis d'enrichir progressivement ses fonctionnalités en fonction des retours des utilisateurs.

Cette démarche a permis de réduire les risques, d'améliorer l'adéquation aux besoins métiers et de favoriser une validation continue des choix techniques.

Les principales étapes ont été les suivantes :
\begin{itemize}
  \item \textbf{Cadrage fonctionnel} : recueil des besoins auprès des parties prenantes (secrétariats, équipe technique), identification des contraintes métier et définition des cas d'usage prioritaires ;
  \item \textbf{Prototypage initial} : mise en place d'un premier pipeline minimal (ingestion $\rightarrow$ OCR $\rightarrow$ stockage) afin de valider les choix techniques ;
  \item \textbf{Itérations successives} : ajout progressif des modules d'extraction, d'assistance intelligente (classification, renommage) et de recherche documentaire ;
  \item \textbf{Intégration des interfaces utilisateur} : développement des interfaces de consultation et de validation, avec prise en compte des retours terrain ;
  \item \textbf{Stabilisation et mise en production} : amélioration de la robustesse, mise en place de l'observabilité (logs, métriques, alertes) et préparation au déploiement en environnement \textit{on-premise}.
\end{itemize}

L'implication des parties prenantes a été un élément clé de la démarche. Des échanges réguliers avec les utilisateurs finaux ont permis d'ajuster les fonctionnalités, notamment sur les règles de classement et les interfaces de validation.

Les principaux risques identifiés concernaient :
\begin{itemize}
  \item la qualité et la variabilité des documents scannés (impact sur l'OCR) ;
  \item l'hétérogénéité des formats documentaires ;
  \item la nécessité de garantir un contrôle humain sur les décisions automatisées.
\end{itemize}

Ces risques ont été atténués par :
\begin{itemize}
  \item une architecture asynchrone basée sur un système de messagerie, permettant de fiabiliser le traitement ;
  \item la mise en place d'un workflow de validation utilisateur (\og proposer puis valider \fg{}) ;
  \item une approche progressive permettant d'ajuster les traitements en fonction des résultats observés.
\end{itemize}

\section{Organisation et planification}
Le projet a été structuré selon une logique d'itérations courtes (\textit{sprints}), avec des livraisons incrémentales permettant de valider rapidement les fonctionnalités développées.

Les principaux livrables ont été organisés comme suit :
\begin{itemize}
  \item (i) ingestion automatique des documents ;
  \item (ii) OCR et stockage ;
  \item (iii) extraction des informations et gestion des métadonnées ;
  \item (iv) assistance au classement et au renommage ;
  \item (v) interface utilisateur et moteur de recherche ;
  \item (vi) observabilité, supervision et durcissement du système.
\end{itemize}

Chaque itération donnait lieu à une démonstration auprès des utilisateurs et de l'équipe technique, permettant de recueillir des retours et d'ajuster les développements suivants. Cette boucle de feedback a contribué à améliorer progressivement la pertinence fonctionnelle et l'ergonomie du système.
